\section{Exact algebraic coefficients: a bounded arithmetic construction}
\label{sec:backend}

This section supplies the exact coefficient field used in the main proof.
Its operations, equality tests, bounds, and polynomial splitting are
primitive recursive on finite codes. The only root-existence theorem reused
is the single-root constructive FTA described in Section~\ref{sec:fta}.
It is used once in the completeness argument for each splitting operation,
not as an oracle in the executed algorithm.

\subsection{Real-algebraic codes and a bounded isolation procedure}

A real-algebraic code is either a rational number, or a tuple
\[
 (p,a,b,A,B,R)
\]
of rational polynomials and rationals, with $R$ a positive integer.
The finite validity conditions require $p$ to be monic and nonconstant,
\[
 Ap+Bp'=1,\qquad
 R>1+\sum_{i<\deg p}\abs{p_i},\qquad
 -R\le a<b\le R,\qquad p(a)p(b)<0.
\]
Write
\[
 M=1+\sum_i\abs{B_i}R^i,\qquad
 L=1+\sum_{i=2}^{\deg p}\binom i2\abs{p_i}R^{i-2},
 \qquad \sigma=\frac1{2ML},
\]
and additionally require $b-a<\sigma$.
Every validity test is a finite rational calculation; in particular
squarefreeness is certified by the displayed B\'ezout identity.

The represented value is specified by rational bisection on $[a,b]$.
At an exact midpoint zero, all subsequent approximations equal that
midpoint. Otherwise retain the half interval with opposite endpoint signs.
After $s$ steps its length is $(b-a)2^{-s}$. Thus the midpoints have an
explicit Cauchy modulus. If
\[
 D=1+\sum_{i=1}^{\deg p}i\abs{p_i}R^{i-1},
\]
the finite identity for differences of powers gives
$\abs{p(x)-p(y)}\le D\abs{x-y}$ on $[-R,R]$.
The opposite endpoint signs then give
$\abs{p(c_s)}\le D(b-a)2^{-s}$ for the midpoint $c_s$.
These finite inequalities establish a root name without appealing to a
completeness axiom or to an already existing real root.

\begin{lemma}[Effective separation and isolation]
\label{lem:backend-isolation}
All real roots of a nonzero rational polynomial can be returned as a finite
list of valid real-algebraic codes, by a primitive recursive algorithm.
The list is complete even for roots initially supplied only by Cauchy names.
\end{lemma}

\begin{proof}
First compute the squarefree part by rational Euclidean division, and
normalize it to a monic polynomial $p$. Constant polynomials have no roots.
The finite squarefree calculation also certifies that the squarefree part
divides the original polynomial and that a power of the squarefree part,
with exponent at most the original degree, is divisible by the original
polynomial. These identities establish equality of the root sets without
using a disjunction principle for a product of arbitrary real numbers.
Extended Euclid gives $Ap+Bp'=1$. Choose an integer $R$ with the extra margin
$R>2+\sum_{i<\deg p}\abs{p_i}$, and form $M,L,\sigma,D$ as above.
The elementary geometric-sum root bound places all roots strictly in
$(-R,R)$.

At a root $\alpha$, the B\'ezout identity gives
\[
 \abs{p'(\alpha)}\ge\frac1M.
\]
For $\alpha,t\in[-R,R]$, the exact polynomial identity
\[
 p(t)-p(\alpha)-(t-\alpha)p'(\alpha)
       =(t-\alpha)^2 E(\alpha,t)
\]
has $\abs{E(\alpha,t)}\le L$. Indeed, the remainder for the monomial $X^i$
is
\[
 (t-\alpha)^2
 \sum_{k=0}^{i-2}(i-1-k)t^k\alpha^{i-2-k},
\]
and its coefficient sum is $\binom i2$.
If $\alpha,\beta$ are roots with $\abs{\alpha-\beta}\le\sigma$, this identity
implies
\[
 \frac{\abs{\alpha-\beta}}M
 \le L\abs{\alpha-\beta}^2
 \le \frac{\abs{\alpha-\beta}}{2M}.
\]
Hence $\alpha=\beta$. This argument does not cancel a possibly zero real.

Choose an integer $K>32ML$ and put $h=1/K<\sigma/16$.
Evaluate $p$ at every grid point $j/K$ in $[-R,R]$.
Return each exact grid zero as a rational code. For each cell with nonzero,
oppositely signed endpoint values, return its bisection code.
Its width is less than $\sigma$, so the preceding estimate certifies
uniqueness. A cell adjacent to a grid zero contains no other root.

For completeness, let $\alpha$ be any named root. If it lies strictly
between nonroot endpoints $l,u$ of a cell, then
\[
 p(t)=(t-\alpha)
       \bigl(p'(\alpha)+(t-\alpha)E(\alpha,t)\bigr).
\]
The expression in parentheses differs from $p'(\alpha)$ by less than
$1/(4M)$ throughout the cell. Consequently the endpoint values have
opposite signs, and local uniqueness identifies $\alpha$ with the
returned bisection root.

The selection of this cell does not assume decidable equality of an
arbitrary named real with a grid point. Compute the positive rational
\[
 \mu=\min_{\substack{g\text{ a grid point}\\p(g)\ne0}}
                         \frac{\abs{p(g)}}D.
\]
The list in this minimum is nonempty, since its two extreme grid points
are not roots. Every root has distance at least $\mu$ from every nonroot
grid point. Obtain a rational approximation $c$ to $\alpha$ with error
less than $\min(\mu,\sigma,h)/8$.
If $\abs{c-g}\le\sigma/2$ for a grid zero $g$, then
$\abs{\alpha-g}<\sigma$, and the local uniqueness estimate proves
$\alpha=g$. Otherwise $\alpha$ is positively separated from every grid
zero as well. The rational value $c$ now determines its cell, whose
endpoints are nonroots. This gives a bounded selection in all cases.

The list has no repetitions: different cells have disjoint interiors,
and a nonroot grid endpoint is positively separated from every root.
Distinct returned roots are apart. Cancelling their positively bounded
distance in the same remainder estimate gives separation at least
$1/(ML)$; the weaker bound $\sigma$ will suffice below.
All searches range over the predetermined $2RK+1$ grid points.
\end{proof}

\subsection{Arithmetic, equality, and quantitative bounds}

Let $\alpha,\beta$ have monic rational annihilators $p,q$.
Finite resultant identities give annihilators for their sum and product:
\[
 \Res_T\bigl(p(T),q(Z-T)\bigr),\qquad
 \Res_T\bigl(p(T),T^{\deg q}q(Z/T)\bigr).
\]
Both are monic in $Z$, of degree $(\deg p)(\deg q)$, and therefore nonzero.
Negation uses $p(-Z)$. For a certified nonzero $\alpha$, remove all powers
of $T$ from $p$, obtaining $p_0$ with $p_0(0)\ne0$; an annihilator for
$1/\alpha$ is
$Z^{\deg p_0}p_0(1/Z)$.
The latter polynomial is nonzero. The annihilation identities follow by
polynomial elimination and do not presuppose a complex splitting algorithm.

Apply Lemma~\ref{lem:backend-isolation} to the chosen annihilator.
Using the input bisection names, approximate the intended arithmetic result
within $h/4$, where $h$ is the new isolating grid's mesh.
Each candidate's rational midpoint is within $h/2$ of its root, while
distinct candidates are separated by more than $16h$.
Exactly one midpoint therefore lies within $h$ of the computed
approximation. Finite rational comparisons select its code.
For products, input root bounds give the required input precision.
For reciprocals, use the explicit positive lower bound established next.

\begin{lemma}[Decidable algebraic zero and sign]
\label{lem:backend-zero}
Zero, equality, and order are primitive recursive for real-algebraic codes.
A nonzero value comes with a positive rational lower bound on its absolute
value.
\end{lemma}

\begin{proof}
Let a named value $\gamma$ have a nonzero rational annihilator
$r(Z)=Z^k q(Z)$ with $q(0)\ne0$, obtained by removing the initial zero
coefficients. Put
\[
 \eta=\min\left(1,
       \frac{\abs{q(0)}}{1+\sum_{i\ge1}\abs{q_i}}\right)>0
\]
and compute a rational approximation $c$ with error at most $\eta/8$.
If $\abs c<\eta/2$, then $\abs\gamma<5\eta/8$, and
\[
 \abs{q(\gamma)-q(0)}
 \le\abs\gamma\sum_{i\ge1}\abs{q_i}
 <\frac58\abs{q(0)}.
\]
Thus $q(\gamma)$ is positively apart from zero. Since
$\gamma^k q(\gamma)=0$, it follows that $\gamma^k=0$ and $\gamma=0$.
This uses the elementary finite power estimate, not a
nonzero-to-apartness principle. If instead $\abs c\ge\eta/2$, then
$\abs\gamma\ge3\eta/8$, and $\gamma$ has the sign of $c$.
Both alternatives are decided by one rational comparison.
Apply this procedure to the difference of two codes to obtain equality
and order. The argument also applies directly to the named arithmetic
result before its output isolating code has been selected.
\end{proof}

Upper bounds are supplied by the root bound $R$.
Complex-algebraic codes are pairs of real-algebraic codes, representing
$a+ib$; their equivalence is decidable equality of both components.
The usual two-coordinate formulas give addition and multiplication.
For a nonzero pair, at least one component has a positive rational
absolute-value lower bound, which also bounds the complex modulus from
below. An upper bound is supplied by the sum of the two component bounds.
Inversion uses
\[
 (a+ib)^{-1}=\frac{a-ib}{a^2+b^2},
\]
whose denominator is certified positive by the same tests.
This constructs the represented field $\A$ and all operations required
in the main argument.

\subsection{A rational annihilator for every root}

Consider a monic polynomial
\[
 H(Y)=Y^n+\sum_{j<n}(a_j+ib_j)Y^j
\]
over the represented field. Choose monic rational annihilators $p_j,q_j$
for its real and imaginary coefficient parts. Form the finite rational
algebra
\[
 T=\Q[U_0,V_0,\ldots,U_{n-1},V_{n-1},J]/
       \bigl(p_j(U_j),q_j(V_j),J^2+1\bigr).
\]
Its explicit monomial basis has dimension
$D_0=2\prod_{j<n}(\deg p_j)(\deg q_j)$.
The coefficient tuple defines an evaluation homomorphism sending
$U_j,V_j,J$ to $a_j,b_j,i$. No primitive-element theorem or compatible
embedding search is involved.

Let $M(Y)$ be the matrix of multiplication on this basis by
\[
 H_T(Y)=Y^n+\sum_{j<n}(U_j+JV_j)Y^j,
 \qquad B_0(Y)=\det M(Y).
\]
Since $M(Y)=Y^nI$ plus lower-degree terms, $B_0$ is monic of degree $nD_0$.
In particular it is nonzero, even though $T$ may have zero divisors.
Apply the adjugate matrix to the coordinate vector of $1$ to obtain an
element $E(Y)\in T[Y]$ satisfying
\[
 H_T(Y)E(Y)=B_0(Y)\,1.
\]
Evaluation at any named root $\beta$ of $H$ proves $B_0(\beta)=0$.
The construction is a finite determinant calculation with rational
coefficients. A nonmonic polynomial is first normalized by its certified
nonzero leading coefficient.

\subsection{Finite candidates and complete splitting}

Write $m=\deg B_0$. Define
\begin{align*}
 R_0(U)&=\Res_Z\bigl(B_0(Z),B_0(2U-Z)\bigr),\\
 T_0(V)&=\Res_Z\bigl(B_0(Z),B_0(Z-2iV)\bigr)\in\Q(i)[V],\\
 S_0(V)&=T_0(V)\overline{T_0}(V)\in\Q[V],
\end{align*}
where the bar conjugates coefficients.
The leading coefficients of $R_0,T_0$ are respectively
$2^{m^2}$ and $(-2i)^{m^2}$, and both have degree $m^2$.
Consequently $R_0$ and $S_0$ are nonzero.
These assertions follow by selecting the top parameter-degree terms in
the resultant determinant.

If $\beta=u+iv$ is a root of $B_0$, its conjugate is also a root because
$B_0$ has rational coefficients. Substitution at $Z=\beta$ in the
resultant identities gives
\[
 R_0(u)=0,\qquad T_0(v)=0,\qquad S_0(v)=0.
\]
Isolate all real roots $r_1,\ldots,r_A$ of $R_0$ and
$s_1,\ldots,s_B$ of $S_0$. The finite Cartesian list
\[
 \{\,r_j+i s_k:1\le j\le A,\ 1\le k\le B\,\}
\]
therefore contains every named complex root of $H$.
Completeness here includes a constructive finite identification of both
coordinates, by Lemma~\ref{lem:backend-isolation}.
Extraneous candidates caused by other algebraic embeddings are harmless.

Start with residual polynomial $H$. Process the candidate list in order.
At a candidate $\beta$, use exact equality to test the residual value.
Whenever it is zero, divide exactly by $Y-\beta$, record the root,
and repeat, for at most $n$ divisions in total.
This produces the exact identity
\[
 H(Y)=\prod_{\text{recorded }\beta}(Y-\beta)\,H_{\rm rem}(Y).
\]
The residual is monic. No processed candidate becomes a root again:
if it vanished after a later division, the preceding polynomial would
also have vanished there. Thus no candidate is a root of $H_{\rm rem}$.

Its degree is a decidable natural number. If that degree were positive,
the single-root constructive FTA of Section~\ref{sec:fta}, applied to
$H_{\rm rem}$, would supply a named root $\gamma$.
Its hypotheses hold: a single fixed primitive recursive evaluator reads
the algebraic coefficient codes, and the leading coefficient is $1$.
The displayed factorization implies $H(\gamma)=0$, so the candidate
construction identifies $\gamma$ with an already processed candidate.
This contradicts the residual tests. Hence $H_{\rm rem}=1$.
The recorded list is a complete splitting with multiplicities.

All executed loops are bounded: Euclidean steps by degree, the rational
grid by $2RK+1$, bisection by its explicit precision, determinants by their
finite dimensions, and synthetic divisions by $n$.
Finite code validity and the loop invariants are rational Boolean tests.
The correctness proof is therefore available to quantifier-free induction
after a precision horizon is fixed. Infinite names abbreviate their
explicit finite estimates. The reused FTA theorem supplies completeness
of this finite algorithm; it does not replace the exact coefficient
backend, and it does not introduce a universal evaluator for arbitrary
primitive recursive programs.
